\documentclass[APA,STIX1COL]{WileyNJD-v2}

\articletype{Special Issue Paper}

\received{}
\revised{}
\accepted{}

\raggedbottom

\begin{document}

\title{Extreme Maximum Temperature Events in Mexico: A Weibull Distribution Approach with Adaptive Plotting Positions}

\author[1]{Sazcha Olivera}

\authormark{OLIVERA \textsc{et al}}

\address[1]{\orgdiv{Department of Design Theory and Processes, Communication Sciences and Design}, 
\orgname{Universidad Aut\'onoma Metropolitana--Cuajimalpa},
\orgaddress{\state{Mexico City}, \country{Mexico}}}

\corres{Sazcha Olivera, Department of Design Theory and Processes, Communication Sciences and Design, Universidad Aut\'onoma
Metropolitana--Cuajimalpa, 05370 Mexico City, Mexico.
\email{satzcha@msn.com}}

%% ─────────────────────────────────────────────────────────────
\abstract[Abstract]{
An adaptive methodology is proposed for modeling the right tail of
daily maximum temperature exceedances using a two-parameter Weibull
mixture model, extending the weighted least-squares framework of
\citet{olivera2019} along three dimensions:
(i)~the reference threshold is estimated as the modal temperature
via kernel density estimation (KDE);
(ii)~the number of bins in the grouped histogram is determined
by the Freedman--Diaconis rule; and
(iii)~the parameter~$a$ of the symmetric plotting-position
family $p_i(a)=(i-a)/(n+1-2a)$ is optimized by adapting the
plotting position to the data, an operation theoretically justified
for the subfamily $0<k<1$.
The methodology is applied to daily maximum temperature records
from five cities in Mexico and the southwestern United States
across three twenty-year periods.
Results show that the modal threshold~$\mu^*$ exhibits an upward
trend in most of the cities analyzed, while the shape
parameter~$k$ remains comparatively stable within each
region, indicating that changes in thermal extremes are expressed
primarily as shifts in the typical temperature level.
}

\keywords{Weibull distribution; plotting positions; parameter~$a$;
kernel density estimation; extreme temperatures; Weibull tail;
class construction; Freedman--Diaconis; Bergman weights.}

\jnlcitation{\cname{\author{Olivera S.}}
             (\cyear{2026}),
             \ctitle{Extreme Maximum Temperature Events in Mexico: A Weibull
             Distribution Approach with Adaptive Plotting
             Positions},
             \cjournal{[Journal]}, \cvol{00:1--XX}.}

\maketitle

%% ════════════════════════════════════════════════════════════
\section{Introduction}\label{sec:intro}
%% ════════════════════════════════════════════════════════════

The Weibull distribution occupies a prominent place in applied
statistics because it combines a relatively simple formulation
with notable flexibility of shape. Since its earliest applications
to the analysis of material strength \citep{weibull1951}, its use
has extended to hydrology, climatology, reliability engineering,
and risk analysis. In all of these fields, interest lies not in
the center of the distribution but in its right tail: the
probability that the variable exceeds high thresholds, rarely
observed but with significant consequences.

When the Weibull distribution is estimated via the probability
plot, the ordered observations do not, by themselves, carry their
cumulative probabilities. Each $x_{(i)}$ must be assigned a
non-exceedance probability, known as a \emph{plotting position}~$p_i$.
The choice of $p_i$ is not innocuous: it determines the coordinates
of the linear regression on the probability plot and, consequently,
the estimates of the shape~$k$ and scale~$\lambda$ parameters.

A family that encompasses many of the classical rules is
\begin{equation}\label{eq:familia_intro}
p_i(a) = \frac{i - a}{n + 1 - 2a},\qquad 0 \leq a < \tfrac{1}{2}.
\end{equation}
The value $a=0$ generates the Weibull rule~\citep{weibull1951};
$a\approx 0.30$ corresponds to the Benard and
Bos-Levenbach approximation~\citeyearpar{benard1953};
$a=0.375$ yields the Blom formula~\citeyearpar{blom1958};
$a\approx 0.44$ gives Gringorten~\citeyearpar{gringorten1963};
and the limit $a\to 0.5$ reproduces the Hazen position.
Although these rules differ little numerically, their impact on
estimated extreme quantiles can be substantial, particularly
when the sample size is moderate and the distribution has a
heavy tail.

There is at the same time a well-established theoretical basis
for plotting positions. If $X_{(i)}$ is the $i$th order statistic
of a continuous sample, we define $U_{(i)}=F(X_{(i)})$ such that

\begin{equation*}
    U_{(i)}\sim\mathrm{Beta}(i,\,n+1-i)
\end{equation*}

and

\begin{equation}
    \mathbb{E}[U_{(i)}]=\frac{i}{n+1},
\end{equation}

\medskip

which explains the Weibull rule \citep{makkonen2014}. However, the
linearized fit uses the double-logarithmic transformation
$g(p)=\ln[-\ln(1-p)]$, which is strictly convex. Consequently,
Jensen's inequality implies that $g(\mathbb{E}[U_{(i)}])\neq
\mathbb{E}[g(U_{(j)})]$, and the position that is exact on the
probability scale does not minimize the bias on the transformed
scale. The magnitude of this gap depends on the curvature of~$g$
in the relevant probability range and therefore on the shape
parameter~$k$ of the fitted model.

\citet{goda1988,goda2010} determined through Monte Carlo simulation
that the optimal plotting position for Weibull extremes depends
explicitly on~$k$, providing calibrated coefficients for extreme
wave analysis. Kim et~al. \citeyearpar{kim2012,kim2015} extended
similar ideas to the generalized logistic distribution. Despite
these precedents, common practice in climatology and hydrology
continues to fix~$a$ in advance or to select the position that
produces the largest~$R^2$ within the sample being analyzed, which
introduces selection optimism without any guarantee of
generalization.

The immediate starting point for this investigation is
\citet{olivera2019}. Their procedure compared plotting positions
for the analysis of precipitation exceedances and used the weight
function derived by Bergman \citeyearpar{bergman1986} to weight
the linearized fit:

\begin{equation}\label{eq:bergman_intro}
w_i = (1 - p_i)\,[\ln(1-p_i)]^2.
\end{equation}

\medskip

This also reveals a characteristic of the weight that is relevant
for extreme value analysis: $w_i$ reaches its maximum at
$p\approx 0.8647$ and tends to zero as $p\to 1$, so it does not
assign the greatest weight to the sample maximum.
\citet{olivera2019} showed that the choice of plotting position
can alter conclusions about trends in extreme quantiles; however,
they treated~$a$ as fixed across all samples.

The present study proposes estimating a rule
$\hat{a}=g(n,\tilde{k})$ within the symmetric
family~\eqref{eq:familia_intro}, where $n$ is the number of
classes and $\tilde{k}$ is a preliminary estimate of the shape
parameter obtained without plotting positions. The primary domain
is the Weibull with $0<k<1$: in this interval the distribution
exhibits a stretched-exponential tail that decays more slowly than
an exponential but does not constitute a power-law (Pareto) tail.
This Weibull tail regime \citep{gardes2005,diebolt2008} is the
empirically relevant one for maximum temperature exceedances in
continental and semi-arid climates, where intense heat episodes
are frequent but extreme heat has a physical upper bound that
rules out power-law behavior.

The methodology integrates, in addition to the optimization
of~$a$, three data pre-processing contributions that address
limitations of earlier implementations:

\textbf{1. Modal threshold via KDE.}
The threshold separating normal days from exceedances is
estimated as the modal temperature using the kernel density
estimator \citep{rosenblatt1956,parzen1962}. This reflects the
most frequent temperature and allows the proportion of exceedances
to be an estimated quantity rather than one fixed in advance.

\textbf{2. Number of classes via Freedman--Diaconis.}
The bin width used in the grouped regression is determined by the
Freedman--Diaconis rule~\citeyearpar{freedman1981}, which is
asymptotically optimal under weaker conditions than Sturges' rule
and adapts the number of classes~$n$ to the size and spread of
the data for each city and period.

\textbf{3. Mixture model.}
The exceedance variable $T_i=\max(X_i-\mu^*,0)$ is not purely
continuous: it carries a point mass at $T=0$ for all days with
$X_i\leq\mu^*$ and a continuous Weibull component for $T>0$.
The explicit specification of the mixture distribution

\begin{equation}\label{eq:mezcla_intro}
F_{\mathrm{mix}}(t) = p_0 + (1-p_0)\,F(t;\,k,\lambda),
\qquad t>0,
\end{equation}

is necessary so that the empirical and fitted distribution
functions are compared on the same scale. Without this step, the
empirical CDF starts at~$p_0\approx 0.5$ while the Weibull CDF
starts at~0, producing artificially large residuals.

The central research question is: \textit{Can changes in the right
tail of the maximum temperature distribution be detected and
quantified over multiple decades using a Weibull distribution
fitted to threshold exceedances?}
From this question two testable hypotheses are stated:

\begin{enumerate}
\item[\textbf{H1.}] \textbf{Stability of distributional form.}
  The shape parameter~$k$ of the Weibull distribution fitted to
  maximum temperature exceedances does not vary significantly
  across the three periods analyzed within each city. Changes in
  thermal extremes are reflected primarily in shifts of the modal
  threshold~$\mu^*$ and the scale~$\lambda$, not in changes in
  distributional shape.

\item[\textbf{H2.}] \textbf{Displacement of extreme levels.}
  The modal threshold~$\mu^*$ exhibits a positive trend across
  periods in most of the cities analyzed. As a consequence, even
  when distributional shape remains stable, the probabilities of
  exceeding fixed critical thresholds increase across decades,
  with implications for urban cooling energy demand.
\end{enumerate}

The proposed methodology allows both hypotheses to be tested
through comparison of the estimated parameters~$(k,\lambda,\mu^*)$
and the exceedance probabilities computed for climatically
relevant thresholds.

%% ════════════════════════════════════════════════════════════
\section{The Weibull Distribution and the Tail Regime}
\label{sec:weibull}
%% ════════════════════════════════════════════════════════════

To avoid the ambiguity that frequently arises in the literature
from the interchange of the symbols~$\alpha$ and~$\beta$ across
different traditions, the following notation is adopted: $k>0$
for the shape parameter and $\lambda>0$ for the scale parameter
\citep[see also][]{makkonen2014}. The cumulative distribution
function (CDF) and survival function of the two-parameter Weibull
distribution are
\begin{equation}\label{eq:fda}
F(x;\,k,\lambda)
  = 1 - \exp\!\left\{-\left(\frac{x}{\lambda}\right)^{\!k}\right\},
  \qquad x > 0,
\end{equation}
\begin{equation}\label{eq:supervivencia}
S(x;\,k,\lambda)
  = \exp\!\left\{-\left(\frac{x}{\lambda}\right)^{\!k}\right\},
\end{equation}
with hazard rate
\begin{equation}\label{eq:hazard}
h(x) = \frac{k}{\lambda}\left(\frac{x}{\lambda}\right)^{k-1}.
\end{equation}

\medskip
This function is decreasing for $0<k<1$, constant for $k=1$
(the exponential distribution), and increasing for $k>1$.

Following \citet{gardes2005} and \citet{diebolt2008}, a
distribution is said to belong to the \emph{Weibull tail class}
if its survival function satisfies
\begin{equation}\label{eq:cola_weibull}
-\ln S(x) = x^\tau\,\ell(x),\qquad \tau>0,
\end{equation}

\medskip

where $\ell$ is a slowly varying function at infinity. The
standard Weibull is the canonical member of this class with
$\tau=k$ and $\ell\equiv\lambda^{-k}$. For $0<k<1$, the tail is
\emph{heavier than exponential} but \emph{lighter than any
power law} (Pareto type); this intermediate regime rules out
both heavy-tailed Pareto distributions and the pure exponential
($k=1$). For $k>1$, tails decay rapidly; under temperature
exceedances above the modal threshold~$\mu^*$, this reflects
that very large events are infrequent once the typical temperature
is exceeded.

The regime of~$k$ also has a direct implication for the choice
of plotting position. \citet{goda1988,goda2010} showed through
Monte Carlo simulation that the optimal position for Weibull
extremes depends on~$k$ and that the dependence is most
pronounced when $k<1$, which motivates the adaptive approach
proposed here.

%% ════════════════════════════════════════════════════════════
\section{Modal Threshold and the Mixture Model}
\label{sec:threshold}
%% ════════════════════════════════════════════════════════════

\subsection{Motivation}

Let $\{X_i\}_{i=1}^{N}$ denote the series of observed daily
maximum temperatures in a given city and period. To construct the
exceedance variable, a reference value~$\mu^*$ representing the
``typical'' temperature must be fixed: days with $X_i\leq\mu^*$
will be classified as normal and days with $X_i>\mu^*$ as
exceedances. The choice of~$\mu^*$ conditions what proportion of
the observations contributes to tail estimation and, therefore,
the parameters of the fitted distribution.

The most widely used alternative in the literature is the sample
median~$\hat{m}$. Its main limitation is that, by construction,
it forces approximately 50\% of days to be classified as
exceedances, regardless of the actual asymmetry of the
distribution. In cities with an extended right tail---such as
Chihuahua or Phoenix, where extreme heat episodes shift the
50th percentile upward---this causes moderately hot days to be
classified as exceedances, distorting the identification of
genuinely atypical or extreme events.

Instead, the \emph{mode} of the distribution is proposed as
a threshold~$\mu^*$ reflecting the most frequently occurring
temperature. The mode imposes no prior fraction of exceedances;
the value $\hat{p}_0=P(T=0)$ is determined by the data
themselves. However, the empirical mode is unstable: small
changes in the rounding criterion can shift it by several tenths
of a degree, and with thousands of records it may not be unique.
The frequency histogram avoids the rounding problem but introduces
arbitrariness in the choice of interval boundaries and produces a
step-function curve, which is inadequate for locating a
continuous maximum.

\emph{Kernel density estimation} (KDE) resolves both
difficulties: it smooths the empirical distribution, producing a
differentiable curve whose maximum constitutes the proposed
estimator~$\mu^*$.

\subsection{The Kernel Density Estimator}\label{subsec:kde}

Let $K:\mathbb{R}\to[0,\infty)$ be a kernel function and $h>0$
the bandwidth controlling the width of each bell. The KDE at a
generic point~$x$ is defined as \citep{rosenblatt1956,parzen1962}:

\begin{equation}\label{eq:kde}
\hat{f}_h(x)
  = \frac{1}{Nh}\sum_{i=1}^{N}
    K\!\left(\frac{x - X_i}{h}\right).
\end{equation}

\medskip

The standard Gaussian kernel is adopted:
\begin{equation}\label{eq:kde_gauss}
K(u) = \frac{1}{\sqrt{2\pi}}\,e^{-u^2/2},
\end{equation}

\medskip

so that each observation contributes a normal density with mean
$X_i$ and standard deviation~$h$. The estimator is therefore the
equally-weighted mixture of $N$ normal distributions, one
centered at each observation.

The Gaussian kernel is standard in practice because it produces
everywhere-differentiable curves and imposes no compact support,
making it computationally convenient and statistically preferable
when the underlying density is continuous and free of abrupt
discontinuities, as is the case with temperature.

\subsubsection{Scott's Rule and Modal Threshold Estimation}

The parameter~$h$ governs the smoothness of the curve. Too
small a value reproduces each observation individually, while
too large a value flattens the curve until the distributional
structure is lost. An optimal value balances bias and variance.

Formally, the mean integrated squared error has the asymptotic
decomposition

\begin{equation}\label{eq:mise}
\mathrm{MISE}(h)
  \approx \underbrace{\frac{h^4}{4}
            \int[f''(x)]^2\,dx}_{\mathrm{bias}^2}
  + \underbrace{\frac{1}{Nh}
            \int K^2(u)\,du}_{\mathrm{variance}}.
\end{equation}

\medskip

Minimizing with respect to~$h$ shows that the optimum satisfies
$h^*\propto N^{-1/5}$, a universal result that depends on neither
the kernel nor the unknown density. \citet{scott1992} evaluated
the integral $\int[f''(x)]^2\,dx$ under the assumption of a
Gaussian underlying density and Gaussian kernel, obtaining

\begin{equation}\label{eq:scott}
h_{\mathrm{Scott}} = 1.059\,\hat{\sigma}\,N^{-1/5},
\end{equation}

\medskip

where $\hat{\sigma}$ is the sample standard deviation.

Scott's rule is the standard choice when the underlying density
is approximately unimodal \citep{silverman1986}; its assumptions
are acceptable for the maximum temperature distribution in the
periods analyzed.

\medskip

With the estimator~$\hat{f}_h$ available, the mode of the
distribution is estimated as the argument that maximizes the
smoothed density:

\begin{equation}\label{eq:mustar}
\mu^* = \underset{x\in\mathcal{G}}{\arg\max}\;\hat{f}_h(x),
\end{equation}

where $\mathcal{G}$ is a uniform grid of $J=2{,}000$ points in
$[\min X_i, \max X_i]$. The notation $\arg\max$ denotes the
point~$x$ at which $\hat{f}_h$ attains its maximum value.
With $J=2{,}000$ points over a typical range of
$40\,{}^\circ\mathrm{C}$, the resolution is
$\Delta x=0.02\,{}^\circ\mathrm{C}$, five times finer than the
typical precision of thermometers, so discretization errors are
negligible.

\subsection{The Exceedance Variable}\label{subsec:excedencia}

From~$\mu^*$ the exceedance variable is defined as

\begin{equation}\label{eq:Ti}
T_i = \max(X_i - \mu^*,\;0),\qquad i=1,\ldots,N.
\end{equation}

\medskip

The variable~$T_i$ measures by how many degrees the temperature
on day~$i$ exceeds the modal reference value. For days with
$X_i\leq\mu^*$, $T_i=0$. For days with $X_i>\mu^*$, the
exceedance is the positive difference $T_i=X_i-\mu^*$, which can
take any continuous value in $(0,\infty)$. This mixture of a
point mass and a continuous component requires a special
distributional model.

When $T>0$, the conditional distribution of~$T$ given~$T>0$ is
modeled as Weibull with parameters~$(k,\lambda)$:

\begin{equation}\label{eq:dist_cond}
P(T\leq t \mid T>0) = F(t;\,k,\lambda),\qquad t>0.
\end{equation}

The parameter~$k$ in equation~\eqref{eq:dist_cond} characterizes
the shape of the right tail of the original temperature
distribution above~$\mu^*$. Accurately estimating~$k$
and~$\lambda$ from grouped data is the central inferential problem
addressed by the methodology in Section~\ref{sec:estimation}.

\subsection{The Mixture Model}

Standard Weibull distributions are defined over $(0,\infty)$ and
are absolutely continuous: they assign zero probability to any
isolated point. Fitting a Weibull directly to~$T_i$ without
treating the point mass at zero produces an inconsistency, since
the theoretical distribution says $P(T=0)=0$ while the data show
$P(T=0)\approx 0.5$. The fitting procedure attempts to compensate
for this discrepancy by distorting the distribution parameters,
producing biased estimates of~$k$ and~$\lambda$.

Taking the law of total probability as the starting point, for
any value $t>0$ the probability that $T\leq t$ can be decomposed
by conditioning on the two mutually exclusive events $\{T=0\}$
and $\{T>0\}$:

\begin{align}
P(T\leq t)
  &= P(T\leq t\mid T=0)\,P(T=0)
   + P(T\leq t\mid T>0)\,P(T>0) \notag \\
  &= 1\cdot p_0 + P(T\leq t\mid T>0)\cdot (1-p_0),
    \label{eq:pt_total}
\end{align}

\medskip

where the first term equals~1 because if $T=0$ and $t>0$ then
$T\leq t$ always holds, and $p_0=P(T=0)$ has been defined. The
conditional distribution $P(T\leq t\mid T>0)$ is what is modeled
by the two-parameter Weibull: $P(T\leq t\mid T>0)=F(t;\,k,\lambda)$.
Substituting into~\eqref{eq:pt_total} yields the \emph{mixture
model}:

\begin{equation}\label{eq:mezcla}
F_{\mathrm{mix}}(t;\,k,\lambda)
  = \begin{cases}
    p_0 & t=0,\\[6pt]
    p_0 + (1-p_0)\,F(t;\,k,\lambda) & t>0.
    \end{cases}
\end{equation}

\medskip

The Weibull parameters are now interpreted exclusively over the
days that are positive exceedances: $k$ characterizes how quickly
the probability of exceeding a given threshold decreases within
the subpopulation of hot days, and $\lambda$ establishes the
typical scale of those exceedances.

\medskip

The point-mass probability is estimated directly from the data:

\begin{equation}\label{eq:p0}
\hat{p}_0
  = \frac{\#\{i : X_i\leq\mu^*\}}{N}.
\end{equation}

\medskip

Unlike the case with the median as threshold, where
$\hat{p}_0=0.5$ exactly by construction, with the modal
threshold $\hat{p}_0$ is an estimated quantity that depends on
the actual shape of the distribution. In cities with less
skewed distributions, $\hat{p}_0$ may approach 0.5; in climates
with a pronounced hot tail, it can depart substantially, as shown
in Table~\ref{tab:p0}.

\begin{table}[ht]
\centering
\begin{tabular}{lccc}
\toprule
City        & 1940--1968 & 1969--1996 & 1997--2026 \\
\midrule
Mexico City & 0.4796 & 0.4979 & 0.4782 \\
Chihuahua   & 0.7612 & 0.6624 & 0.6828 \\
Guadalajara & 0.4066 & 0.4321 & 0.4693 \\
Monterrey   & 0.7506 & 0.6372 & 0.6759 \\
Phoenix     & 0.8199 & 0.8180 & 0.8190 \\
\bottomrule
\end{tabular}
\caption{Empirical zero-class probability $\hat{p}_0=P(T=0)$
by city and period.}
\label{tab:p0}
\end{table}

%% ════════════════════════════════════════════════════════════
\section{Class Construction and the Role of~$n$}
\label{sec:classes}
%% ════════════════════════════════════════════════════════════

The WLS regression of Section~\ref{sec:estimation} operates on
grouped data: the positive exceedances $\{T_i:T_i>0\}$ are
divided into $n$~classes, each characterized by a representative
value~$t_j$ and a frequency~$f_j$. The number of classes~$n$ is
an implicit parameter of the entire procedure: it affects the
plotting positions assigned to each class, the weights used in
the regression, and ultimately the estimated parameters~$(k,\lambda)$
and the goodness of fit.

In the temperature context, with data spanning several decades
and five cities with different distributions, a universal
fixed~$n$ would be inappropriate. It is therefore treated as a
data-determined quantity.

Two types of class are distinguished.

The first class, containing all observations with $T_i=0$,
corresponds to the probability mass accumulated by all days on
which the temperature does not exceed the modal threshold~$\mu^*$.
Its empirical weight is $\hat{p}_0=P(T=0)=P(X\leq\mu^*)$. Its
representative value is set to
$t_1=\varepsilon=0.01\,{}^\circ\mathrm{C}$ to allow the
logarithmic transformation required by the Weibull linearization.

This class is introduced to model explicitly the mixed structure
of the variable $T=\max(X-\mu^*,0)$: the mixture distribution
$F_{\text{mix}}$ requires that the point mass at zero be separated
from the continuous Weibull component. The zero class does not
enter the weighted linear regression (WLR) that estimates~$k$
and~$\lambda$; its sole purpose is to ensure that the empirical
and fitted CDFs are compared on the same scale.

\medskip

The second type---positive classes---consists of the half-open
intervals $(0,h_{\mathrm{FD}}]$,
$(h_{\mathrm{FD}},2h_{\mathrm{FD}}]$, $\ldots$, with
representative values $t_2=h_{\mathrm{FD}}$,
$t_3=2h_{\mathrm{FD}}$, and so on.

\medskip

The bin width $h_{\mathrm{FD}}$ is determined by the
Freedman--Diaconis rule \citep{freedman1981}:

\begin{equation}\label{eq:FD}
h_{\mathrm{FD}}
  = \frac{2\,\mathrm{IQR}(T)}{N_+^{1/3}},
\end{equation}

\medskip

where $\mathrm{IQR}(T)=Q_{0.75}(T)-Q_{0.25}(T)$ is the
interquartile range of all exceedances $T_i\geq 0$ and $N_+$ is
the number of positive exceedances.

\medskip

Starting from the zero class and the positive classes generated
by equation~\eqref{eq:FD}, the construction continues until empty
classes are encountered. Classes with zero frequency are removed
from the tail to obtain the active set of $n$~classes. The last
class---the one that brings the cumulative frequency to exactly
1---is also excluded from the WLS regression, since it would
require evaluating $Y_j=\ln[-\ln(1-p_j)]$ at $p_j=1$, which is
undefined.

%% ════════════════════════════════════════════════════════════
\section{WLS Estimation and Adaptive Optimization of~$a$}
\label{sec:estimation}
%% ════════════════════════════════════════════════════════════

The Weibull CDF admits an exact linearization. Taking logarithms
twice in equation~\eqref{eq:fda}:

\begin{align}
1 - F(t) &= \exp\!\left\{-\!\left(\frac{t}{\lambda}\right)^{\!k}\right\}
  \nonumber\\[4pt]
-\ln[1-F(t)] &= \left(\frac{t}{\lambda}\right)^{\!k}
  \nonumber\\[4pt]
\ln\!\bigl[-\ln(1-F(t))\bigr] &= k\ln t - k\ln\lambda.
  \label{eq:doble_log}
\end{align}

\medskip

Defining the transformed variables for the $j$th positive class
with representative value~$t_j$ and plotting position~$p_j$:

\begin{equation}\label{eq:XY}
X_j = \ln(t_j),\qquad Y_j = \ln[-\ln(1-p_j)],
\end{equation}

\medskip
equation~\eqref{eq:doble_log} reduces to an exact linear
regression:
\begin{equation}\label{eq:lineal}
Y_j = k\,X_j + b,\qquad b = -k\ln\lambda.
\end{equation}

\medskip

The slope of regression~\eqref{eq:lineal} directly estimates~$k$,
and the intercept~$b$ allows recovery of~$\lambda$:

\begin{equation}\label{eq:parametros}
\hat{k} = \frac{S_{XY}}{S_{XX}},\qquad
\hat{\lambda} = \exp\!\left(-\frac{b}{k}\right)
              = \exp\!\left(-\frac{\bar{Y}_w - \hat{k}\,\bar{X}_w}
                                   {\hat{k}}\right),
\end{equation}

\medskip

where the weighted means are

\begin{align}
\bar{X}_w &= \frac{\sum_j w_j X_j}{\sum_j w_j},\qquad
\bar{Y}_w  = \frac{\sum_j w_j Y_j}{\sum_j w_j}
\end{align}

\medskip

and the weighted sums of products are

\begin{align}
S_{XY} &= \sum_j w_j(X_j - \bar{X}_w)(Y_j - \bar{Y}_w),\quad
S_{XX}  = \sum_j w_j(X_j - \bar{X}_w)^2.\label{eq:sumas_w}
\end{align}

\medskip
The regression coefficient of determination is

\begin{equation}\label{eq:R2_wlr}
R^2_{\mathrm{WLR}} = \frac{S_{XY}^2}{S_{XX}\,S_{YY}},\qquad
S_{YY} = \sum_j w_j(Y_j - \bar{Y}_w)^2.
\end{equation}

\medskip

Following \citet{olivera2019} and \citet{bergman1986}, the weight
function is:

\begin{equation}\label{eq:peso_A}
w_j = (1 - p_j)\,[\ln(1-p_j)]^2.
\end{equation}

\subsection{Adaptive Estimation of~$a$}\label{subsec:familia}

The symmetric one-parameter family $a\in[0,\tfrac{1}{2})$
encompasses the classical estimators:
\begin{equation}\label{eq:familia}
p_i(a) = \frac{i-a}{n + 1 - 2a},\qquad i=1,\ldots,n.
\end{equation}

\medskip
The plotting position of the zero class is always fixed at
$p_1^{(0)}=1/(n+1)$, corresponding to $a=0$ (the Weibull rule).
This ensures that the choice of~$a$ for the positive classes does
not propagate a bias through the class that is excluded from the
regression.

For the positive classes, the parameter~$a$ is determined through
numerical optimization, selecting the value that maximizes the
goodness of fit of the Weibull mixture model. The resulting
parameter is subsequently used in the estimation of the shape and
scale parameters.

\subsection{Operational Algorithm}\label{subsec:algoritmo}

The proposed procedure for computing the Weibull
parameters~$k$ and~$\lambda$ is summarized below.

\begin{enumerate}
\item Estimate~$\mu^*$ via KDE
  (equations~\eqref{eq:kde}--\eqref{eq:mustar})
  and compute $T_i=\max(X_i-\mu^*,0)$.
\item Determine~$h_{\mathrm{FD}}$ (equation~\eqref{eq:FD}) and
  construct the zero class and the positive classes.
\item Remove trailing empty classes and identify~$n$.
\item Obtain preliminary estimates~$(\tilde{k},\tilde{\lambda})$
  by maximum likelihood.
\item Refine the parameters through joint numerical optimization,
  obtaining~$(k,\lambda)$.
\end{enumerate}

%% ════════════════════════════════════════════════════════════
\section{Study Area and Data}
%% ════════════════════════════════════════════════════════════

Daily maximum temperature records were obtained for five cities
with distinct climates within North America: Mexico City (CDMX),
Guadalajara, Monterrey, Chihuahua, and Phoenix (AZ, USA). The
first four are located at different elevations and latitudes
within Mexico; Phoenix represents an arid desert climate at low
elevation in the southwestern United States.

Daily maximum temperature data were obtained from the \emph{ERA5
hourly data on single levels from 1940 to present}
\citep{hersbach2023era5}, a global climate reanalysis produced
by the European Centre for Medium-Range Weather Forecasts (ECMWF)
and distributed through the Climate Data Store of the Copernicus
Climate Change Service (C3S) \citep{hersbach2020era5}. The data
cover the three analysis periods: 1940--1968, 1969--1996, and
1997--2026. All five cities have complete records for all three
periods. The number of observations per city and period varies
with the number of days per year.

Figure~\ref{fig:kde} shows the KDE-estimated density of daily
maximum temperature for each city and period, together with the
modal threshold~$\mu^*$ identified in each case. The temperature
distribution varies considerably across cities, and within each
city the position of the density peak shows an upward trend across
periods in most cases.

\begin{figure}[ht]
    \centering

    \begin{minipage}{0.33\textwidth}
        \centering
        \includegraphics[width=\textwidth]{images/fig1_kde_cdmx_v2.png}
    \end{minipage}
    \hfill
    \begin{minipage}{0.33\textwidth}
        \centering
        \includegraphics[width=\textwidth]{images/fig1_kde_chihuahua_v2.png}
    \end{minipage}
    \hfill
    \begin{minipage}{0.33\textwidth}
        \centering
        \includegraphics[width=\textwidth]{images/fig1_kde_guadalajara_v2.png}
    \end{minipage}

    \vspace{0.5cm}

    \begin{minipage}{0.33\textwidth}
        \centering
        \includegraphics[width=\textwidth]{images/fig1_kde_monterrey_v2.png}
    \end{minipage}
    \begin{minipage}{0.33\textwidth}
        \centering
        \includegraphics[width=\textwidth]{images/fig1_kde_phoenix_v2.png}
    \end{minipage}

    \caption{KDE-estimated density of daily maximum temperature
    for each city and the three periods analyzed. The vertical
    dashed line indicates the modal threshold~$\mu^*$ estimated
    for each period.}
    \label{fig:kde}
\end{figure}

%% ════════════════════════════════════════════════════════════
\section{Results}\label{sec:results}
%% ════════════════════════════════════════════════════════════

The Weibull mixture model fitted by the procedure described in
Section~\ref{sec:estimation} achieves an excellent fit across
all city--period combinations. The coefficient of determination
$R^2$ exceeds 0.99 in every case. The Kolmogorov--Smirnov (KS)
statistic computed on the individual observations $T_i>0$ ranges
from 0.015 to 0.048; in most cases the associated $p$-value is
below 0.05, which reflects the high power of the test at the
sample sizes involved rather than a meaningful lack of fit.

\begin{table}[ht]
\centering
\begin{tabular}{llrrrrrrr}
\toprule
City & Period & $\hat{\mu}^*$ & $k$ & $\lambda$ & $R^2$ & KS & $p_{\mathrm{KS}}$ \\
\midrule
Mexico City & 1940--1968 & 19.54 & 0.884 & 0.367 & 0.9972 & 0.027 & 0.001 \\
            & 1969--1996 & 20.67 & 0.856 & 0.361 & 0.9993 & 0.015 & 0.191 \\
            & 1997--2026 & 22.03 & 0.880 & 0.335 & 0.9971 & 0.023 & 0.006 \\
\midrule
Chihuahua   & 1940--1968 & 30.62 & 0.727 & 0.286 & 0.9978 & 0.035 & 0.005 \\
            & 1969--1996 & 28.50 & 0.730 & 0.172 & 0.9974 & 0.034 & 0.001 \\
            & 1997--2026 & 31.08 & 0.705 & 0.188 & 0.9970 & 0.045 & 0.000 \\
\midrule
Guadalajara & 1940--1968 & 25.00 & 0.862 & 0.244 & 0.9924 & 0.038 & 0.000 \\
            & 1969--1996 & 25.72 & 0.878 & 0.265 & 0.9956 & 0.029 & 0.000 \\
            & 1997--2026 & 27.29 & 0.869 & 0.290 & 0.9961 & 0.029 & 0.000 \\
\midrule
Monterrey   & 1940--1968 & 32.16 & 0.674 & 0.258 & 0.9941 & 0.048 & 0.000 \\
            & 1969--1996 & 29.78 & 0.666 & 0.175 & 0.9951 & 0.048 & 0.000 \\
            & 1997--2026 & 31.95 & 0.652 & 0.196 & 0.9957 & 0.045 & 0.000 \\
\midrule
Phoenix     & 1940--1968 & 38.65 & 0.712 & 0.269 & 0.9961 & 0.045 & 0.001 \\
            & 1969--1996 & 39.09 & 0.699 & 0.237 & 0.9986 & 0.037 & 0.011 \\
            & 1997--2026 & 39.83 & 0.708 & 0.268 & 0.9975 & 0.035 & 0.017 \\
\bottomrule
\end{tabular}
\caption{Estimated Weibull mixture model parameters by city and period.}
\label{tab:parametros}
\end{table}

A cross-city comparison reveals that Monterrey and Phoenix show
the lowest values of~$k$ ($0.65$--$0.72$), indicating relatively
heavier exceedance tails: once the modal temperature is exceeded,
the probability of further exceedances (\textit{extreme events})
decreases more gradually. Mexico City and Guadalajara show the
highest values of~$k$ ($0.86$--$0.88$), with lighter tails
consistent with their elevation and more temperate climate.

\begin{figure}[ht]
  \centering
  \includegraphics[width=0.5\textwidth]{images/fig6_tendencias_v2.png}
  \caption{Temporal evolution of the modal threshold~$\mu^*$ by
  city and period. An upward trend is apparent in most cities.}
  \label{fig:tendencias}
\end{figure}

In Mexico City, $\mu^*$ increases from $19.54\,{}^\circ\text{C}$
in 1940--1968 to $22.03\,{}^\circ\text{C}$ in 1997--2026, an
increment of $2.49\,{}^\circ\text{C}$ over eight decades---the
largest displacement in the study. Guadalajara shows a similar
pattern, with an increase from $25.00$ to $27.29\,{}^\circ\text{C}$
($\Delta\mu^*=2.29\,{}^\circ\text{C}$). Chihuahua and Monterrey
display non-monotonic trajectories, exhibiting a decline in the
second period followed by recovery in the third.
Table~\ref{tab:excedencias} presents the estimated probabilities
of daily maximum temperature exceeding the thresholds of 30, 35,
and 40\,°C.

\begin{table}[ht]
\centering
\begin{tabular}{llrrr}
\toprule
City & Period &
  $P(X>30\,{}^\circ\text{C})$ &
  $P(X>35\,{}^\circ\text{C})$ &
  $P(X>40\,{}^\circ\text{C})$ \\
\midrule
Mexico City & 1940--1968 &  0.28 & 0.016 & 0.001 \\
            & 1969--1996 &  0.37 & 0.015 & 0.001 \\
            & 1997--2026 &  1.50 & 0.109 & 0.007 \\
\midrule
Chihuahua   & 1940--1968 & 27.48 & 2.70  & 0.048 \\
            & 1969--1996 & 25.03 & 3.64  & 0.260 \\
            & 1997--2026 & 38.38 & 8.64  & 0.489 \\
\midrule
Guadalajara & 1940--1968 & 12.26 & 1.75  & 0.211 \\
            & 1969--1996 & 14.12 & 1.98  & 0.235 \\
            & 1997--2026 & 21.29 & 2.54  & 0.239 \\
\midrule
Monterrey   & 1940--1968 & 37.89 & 7.42  & 0.106 \\
            & 1969--1996 & 35.62 & 4.44  & 0.114 \\
            & 1997--2026 & 45.97 & 10.97 & 0.268 \\
\midrule
Phoenix     & 1940--1968 & 47.95 & 32.29 & 19.42 \\
            & 1969--1996 & 48.68 & 33.87 & 14.81 \\
            & 1997--2026 & 51.96 & 36.60 & 17.71 \\
\bottomrule
\end{tabular}
\caption{Estimated exceedance probabilities (\%) for critical
temperature thresholds by city and period.
$\dagger$~Empirical estimate ($\theta\leq\mu^*$).}
\label{tab:excedencias}
\end{table}

Changes across periods are especially pronounced for Mexico City:
$P(X>30\,{}^\circ\text{C})$ increases from 0.28\% in 1940--1968
to 1.50\% in 1997--2026, driven primarily by the shift of~$\mu^*$
from 19.5 to 22.0\,°C.

For Chihuahua, the increase in $P(X>35\,{}^\circ\text{C})$ is
particularly pronounced: from 2.70\% in 1940--1968 to 8.64\% in
1997--2026, a factor of 3.2. The increase in
$P(X>40\,{}^\circ\text{C})$ is even larger in relative terms:
from 0.048\% to 0.489\% (factor 10.2), although the absolute
probabilities remain small.

In Phoenix---an inherently exceptional case by virtue of its
desert climate---exceedance probabilities above
$40\,{}^\circ\text{C}$ show a non-monotonic pattern: 19.4\%
(1940--1968), 14.8\% (1969--1996), and 17.7\% (1997--2026).
The reduction in the second period may reflect the slight decrease
in~$k$ (from 0.712 to 0.699) that partially offsets the increase
in~$\mu^*$, although it may also reflect natural variability in
the distribution of extremes.

\begin{figure}[ht]
    \centering

    \begin{minipage}{0.33\textwidth}
        \centering
        \includegraphics[width=\textwidth]{images/fig5_excedencia_cdmx_v2.png}
    \end{minipage}
    \hfill
    \begin{minipage}{0.33\textwidth}
        \centering
        \includegraphics[width=\textwidth]{images/fig5_excedencia_chihuahua_v2.png}
    \end{minipage}
    \hfill
    \begin{minipage}{0.33\textwidth}
        \centering
        \includegraphics[width=\textwidth]{images/fig5_excedencia_phoenix_v2.png}
    \end{minipage}
  \caption{Exceedance probability curves $P(X>\theta)$ by city
  and period for the three cities with the largest inter-period
  changes.}
    \label{fig:excedencia}
\end{figure}

\subsection{Hypothesis Assessment}\label{subsec:hipotesis}

The results allow both hypotheses stated in
Section~\ref{sec:intro} to be assessed.

\textbf{H1.} The shape of the exceedance distribution does not
change substantially over time within each city: the
parameter~$k$ varies by less than 3.5\% across periods within
each city, supporting~H1. However, the decreasing trend of~$k$
in Monterrey (from 0.674 to 0.652) warrants attention: a
progressively heavier tail would imply that the most intense
extreme events become relatively more frequent even in the
absence of a modal-level shift, reinforcing the risk in that city.

\textbf{H2.} Mexico City, Guadalajara, and Phoenix show
monotonic increases in~$\mu^*$ across periods. In these cities,
the displacement of the modal threshold is the primary mechanism
explaining the increase in exceedance probabilities. The most
illustrative case is Mexico City: a $2.5\,{}^\circ\text{C}$ gain
in~$\mu^*$, combined with the stability of~$k$, is sufficient
to multiply the probability of exceeding 30\,°C by a factor of
five.

%% ════════════════════════════════════════════════════════════
\section{Conclusions}\label{sec:conclusions}

This paper proposes and applies an adaptive methodology for
modeling the right tail of the daily maximum temperature
distribution using the two-parameter Weibull distribution. The
results obtained for five cities and three periods allow the
research questions stated in the introduction to be answered.

The modal threshold~$\mu^*$, representing the most frequently
occurring maximum temperature, shows an upward trend in most of
the cities analyzed. In Mexico City, the cumulative increment
between the first period (1940--1968) and the third (1997--2026)
is $2.5\,{}^\circ\mathrm{C}$; in Guadalajara,
$2.3\,{}^\circ\mathrm{C}$; in Phoenix,
$1.2\,{}^\circ\mathrm{C}$.

The shape parameter~$k$ remains comparatively stable within each
city across the three periods (with variations below 4\% in all
cases), supporting hypothesis~H1. It is not possible, with the
methodology employed, to rule out statistically significant
changes in shape without a complementary formal inference
analysis. The changes in~$\mu^*$ are the most systematic and of
the greatest magnitude; changes in~$\lambda$ are of smaller
magnitude and less monotonic; the parameter~$k$ is the most
stable.

Do the probabilities of exceeding critical thresholds change?
Yes, and in some cases considerably. The probability of exceeding
$30\,{}^\circ\mathrm{C}$ in Mexico City increases more than
fivefold between the first and third periods. In Chihuahua, the
probability of exceeding $35\,{}^\circ\mathrm{C}$ triples over
the same interval, reflecting the greater sensitivity of
semi-arid and continental climates---with larger summer thermal
amplitude---to shifts in the modal threshold.

The results suggest that, although the form of the extreme
temperature distribution does not change substantially, the
upward displacement of the modal threshold is sufficient to
produce a perceptible increase in the probability of exceeding
critical levels. Since cooling energy demand (especially for
air conditioning) increases nonlinearly with temperature, moderate
increases in~$\mu^*$ can have consequences for urban energy
infrastructure. This effect is especially relevant in rapidly
growing cities such as Mexico City, where air conditioning
penetration has increased significantly in recent decades.

%% ════════════════════════════════════════════════════════════
\bibliography{refs_weibull}

\end{document}
